% v1.0 manuscript draft. This is intentionally separate from preprint.tex until
% the v1.0 freeze has regenerated every headline result from one tagged commit.
% Build from paper/: pdflatex preprint_v1 && bibtex preprint_v1 && pdflatex preprint_v1 && pdflatex preprint_v1
\documentclass[11pt,a4paper]{article}

\usepackage[margin=2.5cm]{geometry}
\usepackage[T1]{fontenc}
\usepackage[utf8]{inputenc}
\usepackage{graphicx}
\usepackage{booktabs}
\usepackage{array}
\usepackage{enumitem}
\usepackage{microtype}
\usepackage[hidelinks]{hyperref}
\usepackage{xurl}
\usepackage{amsmath}
\usepackage{amssymb}
\usepackage{titling}

\setlength{\droptitle}{-3em}
\predate{}\postdate{}

% Freeze-dependent values. Replace from a machine-generated paper_results.json
% produced by the tagged v1.0 commit. Keeping them as macros prevents stale
% numbers from becoming prose.
\newcommand{\VOneAttackScenarios}{\textsc{TBD}}
\newcommand{\VOneAttackCases}{\textsc{TBD}}
\newcommand{\VOneBaselineCases}{\textsc{TBD}}
\newcommand{\VOneLeaksOff}{\textsc{TBD}}
\newcommand{\VOneLeaksOn}{\textsc{TBD}}
\newcommand{\VOneBenignOn}{\textsc{TBD}}
\newcommand{\VOneCommit}{\textsc{TBD}}

% Values currently reproduced by committed evaluation artifacts; regenerate at
% freeze and change here only if the v1.0 result differs.
\newcommand{\PlannerAccepted}{96\%}
\newcommand{\PlannerPrimary}{92--96\%}
\newcommand{\PlannerAnswerableN}{24}
\newcommand{\PlannerUnanswerableN}{6}
\newcommand{\CompositeCells}{4684}
\newcommand{\CompositeExtraSuppressed}{23}
\newcommand{\CompositeModelPoints}{539}
\newcommand{\CompositeFullBase}{60}
\newcommand{\CompositeFullStrong}{55}
\newcommand{\CompositeCoeffOnly}{42}

\title{A Safe-Outputs Gateway for Untrusted AI Analysts\\in Trusted Research Environments}
\author{Alex R. Wade\\\small\texttt{wade@wadelab.net} \;\textperiodcentered\;
  \small\url{https://github.com/wadelab/safe-tre-agent}}
\date{\small v1.0 manuscript draft \;\textperiodcentered\; preprint, not peer reviewed \\
\small all evaluation data synthetic \;\textperiodcentered\; release commit \VOneCommit}

\begin{document}
\maketitle

\begin{abstract}
Trusted Research Environments (TREs) allow analysts to work with sensitive data
while preventing row-level data from leaving a controlled environment. Existing
statistical disclosure-control workflows largely assume a human analyst: humans
formulate analyses, inspect intermediate results, and submit outputs for release
checking. Introducing a capable language-model agent changes this threat model.
The agent may be prompt-injected by data it reads, may generate code or queries
that expose protected records, and may coordinate sequences of individually
admissible analyses into a disclosure attack.

We present \textit{safe-tre-agent}, a prototype architecture for placing an AI
analyst inside a TRE while treating the model itself as untrusted. The model is
not permitted to execute arbitrary code or construct SQL. Its executable
influence is restricted to a typed declarative query specification over an
allowlisted catalogue. Validated requests execute through a read-only query
engine and pass through statistical disclosure controls, session-level
differencing and query-budget controls, and an auditable release decision before
any result crosses the TRE boundary. Richer statistical procedures, including
generalized linear models and analysis of variance, are fitted only from
disclosure-checked aggregate cells, preventing model fitting from reopening
access to row-level records.

The architecture is stated as a normative specification linked to implementation
and tests. We combine adversarial testing, exhaustive enumeration of finite
procedure spaces, property-based testing, bounded model checking in Alloy, and
theorem proving in Lean 4. The Lean model establishes structural properties of
the query boundary, including identifier exclusion, restriction to declared
read-only views, parameter separation, and limits on the granularity of internal
range predicates. Alloy models stateful disclosure rules and checks
session-level properties and known residual cases. Generated correspondence
artifacts connect both formal representations to the running implementation.

Evaluation on synthetic behavioural datasets asks whether attacks that succeed
without the gateway remain possible with it enabled; whether a model capable of
running locally is sufficient to translate natural-language requests into the
restricted procedure language; and how much analytical utility is lost when a
stronger disclosure checker is composed with the gateway. The result is not a
claim that language models can safely analyse arbitrary sensitive data. It is a
narrower systems claim: \textbf{a useful AI planning layer can be made
non-authoritative, with the disclosure property enforced by a small
deterministic interface whose behaviour can be tested, audited, and partially
proved independently of the model.}
\end{abstract}

\section{Introduction}
\label{sec:introduction}

Trusted Research Environments increasingly mediate access to sensitive health,
administrative, behavioural, and social-science data. Their central security
principle is simple: computation moves to the data rather than the data moving
to the researcher, and only outputs that satisfy disclosure-control policy may
leave the environment. Systems such as OpenSAFELY demonstrate that this model
can operate over population-scale health records~\cite{opensafely2020}.
Statistical disclosure control (SDC) provides a mature family of controls for
limiting what apparently aggregate outputs reveal about individuals~\cite{hundepool2012}.

This architecture was designed around human analysts. A language-model agent
alters the security problem in at least three ways.

First, \textbf{the data become an instruction channel}. An agent processing free
text or other uncontrolled fields can receive adversarial instructions
indirectly through the dataset itself. Indirect prompt injection has already
demonstrated that content incorporated into an LLM-integrated application can
alter model behaviour~\cite{greshake2023}. In a TRE this is not merely an
application-integrity problem: a compromised model is operating beside data
that are deliberately unavailable to the user.

Second, \textbf{analysis becomes cheap to compose}. Classical disclosure
controls commonly reason about individual outputs, while reconstruction and
differencing attacks exploit information distributed over many apparently
innocuous queries~\cite{dinur2003}. A human analyst is rate-limited by attention
and effort. An automated agent is not. Once an agent can issue, inspect, and
adaptively refine queries, the session rather than the individual query becomes
a security object.

Third, \textbf{code generation collapses the intended boundary}. If the agent is
allowed to execute arbitrary Python or SQL over row-level records, an output
checker acts only after arbitrary computation has already occurred. Preventing
a model from printing an identifier is then not enough: an adversarial or
merely mistaken program can encode records into aggregates, model parameters,
exception behaviour, or sequences of outputs.

A natural response is to attempt to make the language model itself trustworthy:
improve its system prompt, train it to refuse unsafe analyses, monitor its
reasoning, or employ a second model as a checker. We take the opposite approach.

\subsection{Treat the model as untrusted}
\label{sec:untrusted}

The central design principle of \textit{safe-tre-agent} is:

\begin{quote}
\textbf{No disclosure guarantee depends on the language model behaving
correctly.}
\end{quote}

The planner may be helpful, mistaken, prompt-injected, or actively adversarial.
Its authority is therefore deliberately small. It can choose among operations
represented by a typed declarative query language, but it cannot extend that
language, execute code, construct raw SQL, access arbitrary columns, decide
whether an output is safe, or determine what crosses the TRE boundary.

This turns the model from a trusted analyst into an \emph{untrusted planner
behind a reference monitor}. A natural-language request is translated into a
\texttt{QuerySpec}: a small structured object identifying a declared dataset,
registered statistical procedure, grouping dimensions, and typed filters. The
specification is validated against a fixed catalogue. A valid specification is
compiled deterministically to a parameterised read-only query over declared
views. The resulting aggregate passes through disclosure control,
session-level accounting, and the release boundary. Invalid specifications
execute nothing.

The model is therefore useful without being authoritative. This separation
matters because model quality and model safety become different questions. A
planner that translates a request incorrectly may produce an irrelevant answer,
which is an analytical-quality failure. A planner that attempts to reference an
identifier or issue a forbidden operation produces no execution, which is a
boundary decision. The former can improve as models improve; the latter need
not depend on that improvement.

\subsection{The session is part of the boundary}
\label{sec:session-intro}

Restricting individual queries is not sufficient. An analyst can sometimes infer
a protected value from the difference between two individually releasable
aggregates. Interactive systems must therefore account for what the user has
already learned.

\textit{safe-tre-agent} maintains per-session lineage for released cohorts and
a bounded query budget. New requests are compared with previous releases where
their outputs are commensurable. Potentially revealing differencing pairs are
denied even when both queries would independently satisfy cell-level disclosure
rules.

This creates a second problem: \textbf{a refusal is itself an output}. If the
decision to refuse depends on hidden information, an adaptive analyst can
potentially learn from the allow/deny pattern. Where possible, the session
auditor therefore bases its decisions on disclosure-safe metadata that are
already available to the analyst. This follows the principle of simulatable
auditing: a denial should disclose no information that could not be
reconstructed from already released information~\cite{kenthapadi2005}.

Where deterministic rules cannot provide an adequate quantitative bound, the
limitation is stated rather than hidden. Differential privacy~\cite{dwork2006}
remains the natural extension for controlling cumulative information release
when exact deterministic outputs and fully simulatable decisions are
insufficient.

\subsection{Useful analysis without reopening the row boundary}
\label{sec:cells-intro}

A system supporting only counts and means would be secure but of limited
research value. Allowing the agent to fit arbitrary models directly to rows,
however, would recreate the code-execution problem.

\textit{safe-tre-agent} therefore treats statistical procedures as
\emph{release contracts}. For procedures such as generalized linear models, the
system first constructs the finite table of sufficient aggregate cells required
for fitting. These cells pass through the ordinary disclosure gateway. Only the
finalized, releasable cells reach the fitting procedure. A released model is
consequently a deterministic function of information that could itself have
been released.

This design has two useful consequences. First, adding a statistical procedure
does not automatically enlarge the trusted computing boundary. The procedure
must declare its admissible request space, output contract, disclosure class,
lineage identity, and finite request skeleton before it becomes reachable.
Second, the information cost of a model becomes visible. A Gaussian model, for
example, requires second moments to estimate dispersion. A cell may therefore
be safe enough to release a mean while unsafe to release its sum of squares. In
that case the model must either omit dispersion-dependent statistics or refuse.
The restriction is inconvenient, but it exposes rather than conceals the
disclosure cost of statistical inference.

\subsection{Contributions}
\label{sec:contributions}

This work makes five contributions:

\begin{enumerate}[leftmargin=*]
  \item \textbf{An agent-specific threat model for TRE analysis.} We identify
  the additional disclosure surfaces created when a language-model agent can
  interpret sensitive data, compose queries, and mediate outputs.

  \item \textbf{A constrained execution architecture that treats the model as
  untrusted.} The planner's executable authority is limited to a typed,
  allowlisted analytical language compiled deterministically to parameterised,
  read-only queries. Disclosure decisions are made outside the model.

  \item \textbf{Stateful release controls for agentic analysis.} The system
  composes conventional cell-level disclosure control with per-session lineage,
  differencing controls, budgets, canonical refusals, and tamper-evident audit
  records.

  \item \textbf{Composable statistical procedures that preserve the boundary.}
  Richer analyses are constructed from gateway-finalized cells rather than raw
  rows, allowing statistical capability to be extended without granting the
  model arbitrary execution.

  \item \textbf{An assurance stack tied to the running implementation.}
  Red-team testing, property tests, exhaustive enumeration, Alloy models, and
  Lean proofs jointly test different classes of claim, with generated
  correspondence checks intended to detect drift between code and formal
  artifacts.
\end{enumerate}

The contribution is intentionally narrower than general AI alignment or
autonomous-agent safety. We do not attempt to prove that the model is aligned,
that arbitrary analyses are safe, or that deterministic SDC prevents all
reconstruction attacks. Instead we ask whether \textbf{model capability can be
separated from disclosure authority} strongly enough to make an AI-mediated
TRE inspectable and testable.

\section{Security model}
\label{sec:security}

\subsection{System setting}

We consider a TRE containing one or more sensitive relational datasets and an
AI-assisted analysis service. A researcher outside the data boundary submits a
natural-language analytical request. A language model inside the environment
proposes a structured analysis. Only released statistical outputs may cross the
boundary.

The protected unit is an individual represented in the data. Direct
identifiers, sensitive free text, internal timestamps, and other prohibited
fields are unavailable through the public analytical catalogue.

The system distinguishes four principals:

\begin{itemize}[leftmargin=*]
  \item the \textbf{researcher}, who submits analytical requests and receives
  approved outputs;
  \item the \textbf{planner}, an LLM that translates a natural-language request
  into a proposed structured analysis;
  \item the \textbf{gateway}, deterministic code that validates, executes,
  checks and accounts for proposed analyses; and
  \item the \textbf{operator/output checker}, responsible for deployment
  controls and residual human review.
\end{itemize}

The planner is inside the computational environment but outside the trusted
policy boundary.

\subsection{Adversary}

We assume the planner may behave arbitrarily. It may be incorrectly prompted,
maliciously instructed, prompt-injected through dataset content, or simply poor
at analysis planning.

We additionally permit the user's sequence of requests to be adversarial. A
user may search for combinations of individually safe outputs that reveal
protected information, deliberately trigger refusals, probe category support,
exploit differences between data views, or attempt to encode row-level
information in apparently statistical outputs.

Sensitive data themselves may contain adversarial strings and pathological
numerical values. Security properties must therefore not depend on free-text
content being benign or numerical measures following the synthetic generator's
typical distribution.

We do \emph{not} assume compromise of the TRE host or operating-system
isolation. The production deployment is also assumed to provide an
authenticated restricted channel, and a production model must execute locally:
sending sensitive prompts or data to a remote model endpoint would itself
constitute egress.

\subsection{Security objective}

The primary objective is \emph{controlled disclosure}, rather than model
correctness:

\begin{quote}
No information should cross the release boundary except through a registered
output contract after the applicable disclosure and session controls have been
applied.
\end{quote}

This objective has several components:

\begin{enumerate}[leftmargin=*]
  \item a valid request cannot name a direct identifier or prohibited field;
  \item the planner cannot execute arbitrary code or SQL;
  \item output tables must satisfy the configured cell-level disclosure policy;
  \item sequences of releases are subject to session-level accounting;
  \item information contained in denial behaviour, metadata, and timing is
  considered part of the output surface; and
  \item every request and release decision is attributable and auditable.
\end{enumerate}

The system does not claim information-theoretic privacy. Exact aggregate
releases inevitably reveal information about the underlying data. Deterministic
SDC constrains that release according to policy; it does not provide the
compositional guarantee of differential privacy.

\subsection{Non-goals and trust assumptions}

The following are explicitly outside the v1.0 claim.

\paragraph{Model alignment.} The planner does not need to follow security
instructions.

\paragraph{Arbitrary Python analysis.} A legacy code-writing demonstration is
retained for comparison but is not part of the secure path.

\paragraph{Protection against host compromise.} OS isolation, host security,
and the upstream identity system are deployment assumptions rather than
properties of the analytical gateway.

\paragraph{Cross-user collusion.} Session controls limit adaptive analysis
within the configured accounting scope; they do not constitute a complete
global privacy accountant.

\paragraph{Production certification.} The implementation is a research
prototype using synthetic data. It demonstrates and tests an architecture; it
is not represented as a certified TRE product.

\section{Architecture}
\label{sec:architecture}

The secure path has seven conceptual stages:

\begin{center}
\textbf{request $\rightarrow$ vetting/fidelity $\rightarrow$ planner $\rightarrow$
validation $\rightarrow$ read-only execution $\rightarrow$ disclosure/session
controls $\rightarrow$ release/audit}
\end{center}

The crucial boundary occurs after planning. The planner emits data; it does not
emit executable authority.

\subsection{QuerySpec}
\label{sec:queryspec}

\texttt{QuerySpec} is a declarative representation of the allowed analytical
operation. Its vocabulary is generated from a public catalogue describing
available datasets or views, permitted grouping dimensions, permitted measures,
typed filter operations, registered statistical procedures, and public metadata
required to formulate requests.

Unknown fields are rejected. Columns absent from the appropriate allowlist
cannot be referenced. Internal fields may be available for restricted filtering
without becoming releasable dimensions. Direct identifiers and free text are
excluded from every executable allowlist.

The system also accepts a literal \texttt{QuerySpec} supplied directly by an
analyst. This is useful both operationally and conceptually: it demonstrates
that the language-model planner is optional. Removing the model changes
usability but not the downstream disclosure boundary.

\subsection{Compilation and execution}
\label{sec:execution}

A validated specification is compiled to a restricted read-only query
representation and then to parameterised SQL. The compiler has no
representation for DDL or DML, analyst-selected joins, subqueries, arbitrary
expressions, or user-provided SQL fragments. Filter values are bound as
parameters rather than interpolated into query text.

Execution occurs only over declared public views and under bounded resource
settings. The result returned by the engine is not itself releasable; it is an
internal object supplied to the disclosure gateway.

\subsection{The release gateway}
\label{sec:gateway}

The gateway applies disclosure rules to each candidate output. The v1.0
implementation supports conventional controls including minimum distinct-person
cell size, contribution/dominance rules for additive outputs, influence checks
for correlation, count rounding, primary and complementary suppression, and
optional composition with an external SDC checker.

The purpose of this layer is not to claim a universally correct disclosure
policy. TREs differ in their rules. Instead, the gateway is the policy seam:
the model cannot bypass it, and additional checkers can strengthen cell
decisions without expanding the model's authority.

\subsection{Session accounting}
\label{sec:session}

A release that is safe in isolation may be unsafe in combination with an
earlier release. Each session therefore retains lineage describing previously
released cohorts and the quantities they measured.

When a new result is considered for release, the session auditor asks whether
it is sufficiently close to a previous commensurable release to permit a
differencing attack. The accounting identity follows the underlying population
and declared measured quantity rather than simply the dataset/view name,
because two different views may expose the same per-person quantity.

The system additionally imposes finite query and selection budgets. Accounting
state is included in the audit trail and reconstructed after restart so that
restarting the service cannot reset a user's disclosure history.

\subsection{Refusals are outputs}
\label{sec:refusals}

A policy that says only ``release unsafe outputs less often'' is incomplete.
The pattern of refusals can itself reveal hidden properties such as whether a
cohort contains fewer than a threshold number of people.

\textit{safe-tre-agent} therefore distinguishes decisions that can be justified
from public or previously released information from decisions that depend on
protected data. Data-dependent refusals return a canonical response rather than
detailed internal findings. Detailed execution traces remain available to the
output-checking side of the boundary through the audit record.

This principle extends beyond text. Response latency, exception classes, row
ordering, and suppression choices are observable outputs and are therefore
treated as potential disclosure channels.

\section{Extending the analytical language without extending trust}
\label{sec:procedures}

\subsection{Procedures as registered contracts}

Every analytical operation reachable through the secure path is a registered
procedure. A procedure declares its admissible request schema, compilation and
execution method, disclosure class, output columns, lineage identity, and a
finite request skeleton where the procedure admits finite enumeration. An
unregistered procedure fails closed.

This makes extension explicit. Adding regression is not equivalent to allowing
arbitrary model-fitting code: it requires defining what information the
regression consumes, what it may emit, and how those objects pass through the
disclosure model.

\subsection{Cells-first model fitting}
\label{sec:cells-first}

Generalized linear models illustrate the approach. Instead of fitting a model
to individual observations, the system first constructs a design table of
aggregate cells. Those cells are passed through the ordinary disclosure
gateway. Suppressed cells are not available to the fitter.

The fitting procedure then consumes only the finalized design-cell table. The
released model contains that table alongside the coefficient output, allowing
the fit to be reproduced from released-equivalent information alone.

This provides a useful invariant:

\begin{quote}
If two underlying datasets produce the same finalized input cells, the model
procedure must produce the same released model.
\end{quote}

The fitter therefore does not create a second route from raw records to
released coefficients.

\subsection{Disclosure has an analytical cost}

This architecture makes some limitations unusually explicit. For Gaussian
models, coefficient estimation and uncertainty estimation require different
aggregate information. Means and counts may be releasable where second moments
are not. Consequently the v1.0 procedure can distinguish a complete model
release, a coefficients-only release where dispersion information is withheld,
and refusal where the required design itself cannot be safely represented.

This is preferable to silently computing uncertainty from information the
analyst was not otherwise permitted to obtain. The same principle extends to
analysis of variance and other registered procedures: the statistics released
must be functions of vetted intermediate objects with a stated disclosure
class.

\section{Assurance}
\label{sec:assurance}

No single verification technique matches all of the system's claims. v1.0
therefore deliberately uses several.

\subsection{Normative specification}

The repository maintains a specification of assumptions, requirements,
prohibitions, and non-goals. Each clause has a stable identifier and maps to
enforcing code and tests. This makes the security claim reviewable and prevents
a growing prototype from silently changing what ``safe'' means as new
procedures are added.

\subsection{Adversarial testing}
\label{sec:redteam}

The red-team harness compares execution with the gateway disabled and enabled.
Attack scenarios include direct row extraction, over-granular queries,
adversarial text, query composition, differencing, malformed requests,
model-specific attacks, cross-view attacks, and adversarial numerical fixtures.

The oracle is deliberately independent of the gateway's own findings. A test
passes only when the released information does not satisfy the attack objective;
merely observing that some control fired is insufficient.

At the v1.0 freeze, the paper reports one machine-generated result record from
the tagged commit:

\begin{table}[ht]
\centering
\caption{Adversarial evaluation. Values are populated from the v1.0 release
artifact rather than copied manually into prose.}
\label{tab:redteam}
\begin{tabular}{lrr}
\toprule
 & Gateway disabled & Gateway enabled \\
\midrule
Attack scenarios & \VOneAttackCases & \VOneAttackCases \\
Successful disclosures & \VOneLeaksOff & \VOneLeaksOn \\
Benign baselines released & -- & \VOneBenignOn \\
\bottomrule
\end{tabular}
\end{table}

The disabled condition is not intended to represent a plausible production
configuration. It establishes that the attack fixture contains real disclosures
and prevents a trivially safe test suite in which no attack was capable of
succeeding in the first place.

\subsection{Property and exhaustive testing}

Property-based tests probe numerical and structural edge cases that hand-written
examples are unlikely to contain. Registered finite request skeletons are
additionally enumerated where feasible. Exhaustive enumeration is particularly
useful for checking the closure properties of a deliberately small analytical
language: every admissible structural request can be exercised rather than
sampled.

\subsection{Alloy}
\label{sec:alloy}

Alloy is used where the desired property is naturally expressed as bounded
relational or temporal counterexample search. The models cover release
behaviour for statistical models under possible cell-vetting outcomes,
session-level differencing rules, budget and lineage behaviour across requests
and restarts, and explicit residual cases that the system does not claim to
prevent.

Known residuals are encoded as satisfiable examples rather than omitted. A
formal model that accidentally ceases to exhibit a known weakness is treated as
having drifted from the system.

\subsection{Lean}
\label{sec:lean}

Lean 4 is used for structural properties intended to hold over the full query
language. The v1.0 proofs include properties of the following form:

\begin{itemize}[leftmargin=*]
  \item a valid query cannot reference declared identifiers;
  \item internal-only columns cannot become released grouping dimensions;
  \item compiled queries select from one declared view;
  \item user filter values remain bound parameters rather than becoming SQL
  syntax;
  \item restricted internal range predicates cannot cut more finely than their
  declared public bands;
  \item configured disclosure-policy floors imply minimum structural bounds;
  and
  \item release shaping cannot expose internal witness fields that are absent
  from the release type.
\end{itemize}

These are deliberately modest claims. In particular, proving that the query
compiler excludes identifiers is not a proof that the complete application
provides information-theoretic privacy.

\subsection{Binding formal artifacts to code}

Formal verification is useful only if the model describes the implementation
being released. \textit{safe-tre-agent} therefore generates formal catalogue
and skeleton artifacts from the live procedure registry and database views.
Test cases pin compiler output against the implementation, and correspondence
metadata links modeled counterexamples to executable attack tests. The build
fails when these generated representations drift.

\section{Evaluation}
\label{sec:evaluation}

The evaluation is organised around four research questions.

\subsection{RQ1: Does the boundary stop attacks that succeed when it is absent?}

We replay the adversarial corpus with the disclosure boundary disabled and
enabled, measuring released information rather than the gateway's self-reported
findings. Table~\ref{tab:redteam} is regenerated from the tagged v1.0 commit.
The corresponding machine-readable result record should be archived with the
release and DOI.

\subsection{RQ2: Is the restricted language useful enough for a language model to plan into?}
\label{sec:planner-eval}

A secure analytical language is of little practical use if translating ordinary
research questions into it requires a human specialist. We therefore evaluate
natural-language planning against a corpus containing \PlannerAnswerableN
answerable requests and \PlannerUnanswerableN requests for which the analytical
language deliberately has no legal expression.

A 120B-class open-weight model, used as a stand-in for a model that could
plausibly be hosted locally, achieved \PlannerAccepted accepted-spec accuracy
after the dataset semantics and examples required for planning were explicitly
represented in the prompt. Intended-spec accuracy was \PlannerPrimary across
two runs. These values should be interpreted at corpus resolution: with 24
answerable items, one item changes the percentage by roughly four points.

More importantly, planner refusal was unreliable. Models frequently transformed
an unanswerable request into a nearby answerable one rather than refusing. For
example, a request for per-person values could become a request for an overall
mean, or a request for residuals could become a legal regression with the
residual request silently omitted.

This result supports the architectural separation: \textbf{the model is
adequate as a planner but unsuitable as the authority that determines whether a
request is admissible.}

\subsection{RQ3: What utility is lost by composing stronger disclosure checks?}
\label{sec:acro-eval}

The built-in checker can be composed with an external SDC checker. On the
current synthetic evaluation corpus, adding the ACRO-derived checker suppresses
\CompositeExtraSuppressed additional cells among \CompositeCells candidate
cells.

For Gaussian models, the same composition reduces the number of complete fits
from \CompositeFullBase to \CompositeFullStrong among \CompositeModelPoints
enumerated model points, while leaving the \CompositeCoeffOnly
coefficients-only releases unchanged. These values are regenerated at release
freeze.

This is a small but measurable analytical cost, and demonstrates why
disclosure-control composition should be measured rather than assumed to be
free.

\subsection{RQ4: Can automated multi-step analysis remain behind the same boundary?}
\label{sec:inside-analyst}

As an extension experiment, we allow an automated analyst to formulate a
sequence of ordinary requests and assemble their released results into a typed
evidence dossier.

The analyst receives no privileged row-level interface. Each sub-analysis enters
the same query service under a common session, so query validation, disclosure
checking, lineage, and accounting apply exactly as they do to an externally
submitted request.

Claims in the final dossier must cite released analytical steps. A claim without
supporting released evidence is represented as \texttt{not\_answerable}, not
silently inferred from hidden execution state. This experiment is included to
demonstrate composability of the boundary rather than as a claim that
autonomous scientific analysis has been solved.

\section{Discussion}
\label{sec:discussion}

\subsection{The useful model is an untrusted model}

The most important result is architectural rather than numerical. Increasing
planner capability need not require increasing planner authority.

A stronger model can produce better \texttt{QuerySpec}s, formulate more
appropriate sequences of procedures, and explain results more effectively while
remaining unable to bypass the deterministic release boundary. Conversely, a
weaker or compromised model primarily degrades analytical quality rather than
changing which operations are executable. This is a desirable asymmetry for
high-consequence AI deployment.

\subsection{Safe output is a sequence property}

Agentic analysis makes explicit a weakness that already exists in interactive
statistical systems: the unit of disclosure is not necessarily one table. A
sequence of correct release decisions can be collectively unsafe.

The system therefore treats lineage, accounting, refusal behaviour, and restart
semantics as parts of the same security boundary as minimum cell sizes and
dominance rules. This becomes increasingly important as automated analysts
increase the number and adaptivity of analytical requests.

\subsection{Restriction can improve inspectability}

Restricting the analytical language obviously reduces flexibility. It also buys
something unusual in AI systems: a meaningful fraction of the reachable action
space can be enumerated, tested, and reasoned about.

This does not make the system formally secure in the general sense. It does,
however, make particular strong statements possible that would be implausible
for arbitrary model-written Python. The practical design question is therefore
not simply ``how capable can the agent be?'' but:

\begin{quote}
\textbf{What is the smallest authority surface through which the desired
capability can be expressed?}
\end{quote}

\subsection{Disclosure control and model alignment solve different problems}

A model may be cooperative and still make a disclosure mistake. It may be
adversarial and still be unable to escape a sufficiently constrained execution
language. Conversely, a perfectly constrained execution language cannot ensure
that the model interprets the researcher's question correctly.

We therefore distinguish:

\begin{itemize}[leftmargin=*]
  \item \textbf{semantic quality}: did the planner formulate the analysis the
  researcher intended?;
  \item \textbf{execution safety}: was the requested operation within the
  allowed analytical language?; and
  \item \textbf{disclosure safety}: was the resulting information permitted to
  cross the boundary?
\end{itemize}

Conflating these problems makes each harder to evaluate.

\subsection{Why not rely on a model checker?}

One tempting architecture is to let one model perform the analysis and ask a
second model whether the output is safe. That may be useful as a triage aid, but
it does not provide the same property as the deterministic boundary here. Both
models may share failure modes; both can be prompt-injected; and neither gives a
small executable language over which structural properties can be stated.

The architecture presented here can still use learned systems inside the
boundary, but their outputs are treated as proposals or findings rather than as
release authority.

\section{Limitations}
\label{sec:limitations}

The system remains a research prototype.

First, all supplied datasets are synthetic. The implementation has been
deliberately tested with adversarial data fixtures, but deployment against real
protected data would require independent review and adaptation to the host
TRE's disclosure policy.

Second, deterministic SDC does not provide a complete compositional privacy
guarantee. The session auditor prevents specified classes of differencing
attack but does not replace a formal privacy accountant. Differential privacy
is the natural mechanism for bounding cumulative information release where
exact outputs are not required.

Third, the current analytical language is deliberately narrow. Categorical
GLMs and registered aggregate procedures cover useful research questions but do
not approach the flexibility of arbitrary R, Python, or SQL analysis.

Fourth, formal verification covers selected boundary properties rather than the
complete Python, web, and deployment stack. Lean proves properties of the
modeled query and release structures; Alloy checks bounded models. Neither
should be interpreted as a proof of end-to-end security.

Fifth, production deployment depends on environmental assumptions including host
isolation, a local model runtime, trusted identity infrastructure, and correct
operational key management.

Finally, an AI analyst creates scientific-integrity risks that are orthogonal to
disclosure: inappropriate model choice, confounding, data dredging, selective
reporting, and hallucinated interpretation. Constraining what an agent may
release does not make the scientific inference correct.

\section{Related work}
\label{sec:related}

\subsection{Trusted Research Environments and statistical disclosure control}

TREs and related secure-data environments combine controlled computation with
output checking. OpenSAFELY provides a prominent code-to-data architecture
operating at national health-data scale~\cite{opensafely2020}. Classical SDC
provides a mature set of controls for frequency, magnitude, and dominance
disclosure~\cite{hundepool2012}. ACRO automates parts of this checking workflow
and provides an important reference point for machine-assisted SDC~\cite{acro}.

\textit{safe-tre-agent} differs primarily in the assumed analyst. Rather than
regarding an AI assistant as equivalent to a faster human user, we model it as
an untrusted computational principal capable of generating and composing
requests at machine speed.

\subsection{Query auditing and differential privacy}

Database auditing has long recognised that individually innocuous queries may
reveal protected information in combination. Dinur and Nissim demonstrated
fundamental limits of overly accurate statistical query answering~\cite{dinur2003}.
Simulatable auditing addresses the additional information potentially disclosed
by an auditor's allow/deny decisions~\cite{kenthapadi2005}. Differential privacy
provides the canonical formal framework for bounding cumulative disclosure under
composition~\cite{dwork2006}.

Our deterministic session auditor should be understood as an engineering
control for a restricted analytical setting, not as an alternative privacy
definition.

\subsection{LLM agents and indirect prompt injection}

LLM-integrated applications introduce a further trust problem because untrusted
content may affect agent behaviour through indirect prompt
injection~\cite{greshake2023}. Many agent-safety approaches attempt to detect
malicious instructions or make the model more reliable.

Our design does not require successful detection. A prompt-injected planner is
permitted to propose arbitrary objects; the executable boundary accepts only
objects in the same restricted language available to a benign planner.

\section{Reproducibility and release protocol}
\label{sec:repro}

The v1.0 paper is intended to describe one tagged software state rather than a
moving \texttt{main} branch. The release procedure should therefore:

\begin{enumerate}[leftmargin=*]
  \item freeze the v1.0 commit and record its SHA;
  \item install from the locked environment on a clean runner;
  \item run the default, slow/exhaustive, red-team, analyst-red-team, and formal
  checks;
  \item regenerate all figures and evaluation JSON artifacts;
  \item generate a single \texttt{paper\_results.json} containing every numeric
  value used in the abstract, tables, and figures;
  \item build this manuscript from those artifacts; and
  \item archive the tagged software and paper together.
\end{enumerate}

The release artifact should record toolchain versions and hashes for Lean and
Alloy and preserve the planner-evaluation outputs needed to reproduce the
reported ranges. Remote model endpoints are used only on synthetic data; a real
TRE deployment requires a local model runtime.

\section{Conclusion}
\label{sec:conclusion}

AI assistants could substantially reduce the friction of analysing sensitive
data inside Trusted Research Environments. Granting them the authority
traditionally given to a human analyst, however, also gives an unreliable and
potentially adversarial computational process access to a security boundary
designed around human behaviour.

\textit{safe-tre-agent} explores a different architecture. The language model
is useful but untrusted. It proposes analyses through a restricted declarative
interface. Deterministic code decides what may execute and what may leave.
Session state constrains composition across requests. Richer statistical
procedures consume already-vetted aggregate information rather than reopening
row-level access. The resulting boundary is small enough to support adversarial
testing, exhaustive enumeration, and selected formal proofs.

The broader design principle is simple:

\begin{quote}
\textbf{Capability does not require authority.}
\end{quote}

For AI systems operating over sensitive data, separating the two may be more
robust than attempting to make increasingly capable models themselves carry the
security guarantee.

\bibliographystyle{plain}
\bibliography{references}

\end{document}
